% GENERATED FILE. Edit canonical lessons, then run npm run generate:books.
% canonical-source-hash: fnv1a64:75b0273b100f4deb
% canonical-lessons: SA-C439-pavanah, SA-C439-prabhanjanah, SA-C439-prabalah, SA-C439-andolayati, SA-C439-mandanilah

\chapter{Winds and Storms}
\label{ch:sa-winds-and-storms}
\hlchaptermodality{439}

\begin{chapteropening}
\textbf{By the end of this chapter:} \emph{I can say how the wind blows and what it sets swaying.}

Complete the last lesson of chapter 439: I can say how the wind blows and what it sets swaying.
\end{chapteropening}

\section[\texorpdfstring{\sk{पवनः} (pavanaḥ)}{पवनः (pavanaḥ)}]{\sk{पवनः} (pavanaḥ) — a breeze}

\label{lesson:SA-C439-pavanah}



\begin{quote}
Before the new one: say the Sanskrit for sweat, then the Sanskrit for blazes.
\end{quote}

\subsection*{You'll want to know: \sk{पवनः}}
\textbf{\sk{पवनः}} — \emph{pavanaḥ} — \textquotedblleft{}a breeze\textquotedblright{}.

\begin{quote}
\sk{सायं} \sk{शीतलः} \sk{पवनः} \sk{वहति।}
\end{quote}

\begin{quote}
\emph{sāyaṃ śītalaḥ pavanaḥ vahati.}
\end{quote}

\begin{quote}
In the evening a cool breeze blows.
\end{quote}

A masculine noun.

\begin{cousinweb}
\textbf{\sk{पवनः}} (\emph{pavanaḥ}) comes from the root √pū, \textquotedblleft{}to purify\textquotedblright{}: MW gives \emph{pavana}, \textquotedblleft{}the purifier\textquotedblright{}, as wind, breeze, air; it lives on as Punjabi \emph{pauṇ}, \textquotedblleft{}wind\textquotedblright{} (MW; CDIAL 7978).
\end{cousinweb}

\begin{grammarlens}[title={What you've built}]
One more word for getting around.
\end{grammarlens}

\subsection*{Your turn}
\begin{itemize}
\raggedright
  \item \emph{Say it:} \emph{pavanaḥ}
  \item \emph{Say it:} \emph{pavanaḥ}, once more
  \item \emph{Say it:} \emph{sāyaṃ śītalaḥ pavanaḥ vahati.} — In the evening a cool breeze blows.
  \item \emph{Recall:} say the Sanskrit for sweat, then the Sanskrit for blazes, then say \emph{pavanaḥ} again
\end{itemize}

\begin{tcolorbox}[breakable,colback=teal!4,colframe=teal!35!black,title={Before you move on}]
What does \sk{पवनः} mean? (\textquotedblleft{}A breeze\textquotedblright{}.) Say it once more.
\end{tcolorbox}
\section[\texorpdfstring{\sk{प्रभञ्जनः} (prabhañjanaḥ)}{प्रभञ्जनः (prabhañjanaḥ)}]{\sk{प्रभञ्जनः} (prabhañjanaḥ) — a tempest}

\label{lesson:SA-C439-prabhanjanah}



\begin{quote}
Before the new one: say the Sanskrit for blazes, then the Sanskrit for a breeze.
\end{quote}

\subsection*{You'll want to know: \sk{प्रभञ्जनः}}
\textbf{\sk{प्रभञ्जनः}} — \emph{prabhañjanaḥ} — \textquotedblleft{}a tempest\textquotedblright{}.

\begin{quote}
\sk{रात्रौ} \sk{प्रभञ्जनेन} \sk{वृक्षाः} \sk{पतन्ति।}
\end{quote}

\begin{quote}
\emph{rātrau prabhañjanena vṛkṣāḥ patanti.}
\end{quote}

\begin{quote}
At night trees fall in the tempest.
\end{quote}

A masculine noun.

\begin{cousinweb}
\textbf{\sk{प्रभञ्जनः}} (\emph{prabhañjanaḥ}) is built from \emph{pra} and the root √bhañj, \textquotedblleft{}to break\textquotedblright{}: the breaker; MW gives \emph{prabhañjana} as wind, storm, tempest, hurricane (MW).
\end{cousinweb}

\begin{grammarlens}[title={What you've built}]
One more word for getting around.
\end{grammarlens}

\subsection*{Your turn}
\begin{itemize}
\raggedright
  \item \emph{Say it:} \emph{prabhañjanaḥ}
  \item \emph{Say it:} \emph{prabhañjanaḥ}, once more
  \item \emph{Say it:} \emph{rātrau prabhañjanena vṛkṣāḥ patanti.} — At night trees fall in the tempest.
  \item \emph{Recall:} say the Sanskrit for blazes, then the Sanskrit for a breeze, then say \emph{prabhañjanaḥ} again
\end{itemize}

\begin{tcolorbox}[breakable,colback=teal!4,colframe=teal!35!black,title={Before you move on}]
What does \sk{प्रभञ्जनः} mean? (\textquotedblleft{}A tempest\textquotedblright{}.) Say it once more.
\end{tcolorbox}
\section[\texorpdfstring{\sk{प्रबलः} (prabalaḥ)}{प्रबलः (prabalaḥ)}]{\sk{प्रबलः} (prabalaḥ) — strong; violent}

\label{lesson:SA-C439-prabalah}



\begin{quote}
Before the new one: say the Sanskrit for a breeze, then the Sanskrit for a tempest.
\end{quote}

\subsection*{You'll want to know: \sk{प्रबलः}}
\textbf{\sk{प्रबलः}} — \emph{prabalaḥ} — \textquotedblleft{}strong; violent\textquotedblright{}.

\begin{quote}
\sk{समुद्रे} \sk{प्रबलः} \sk{वायुः} \sk{वहति।}
\end{quote}

\begin{quote}
\emph{samudre prabalaḥ vāyuḥ vahati.}
\end{quote}

\begin{quote}
A violent wind blows over the sea.
\end{quote}

\begin{cousinweb}
\textbf{\sk{प्रबलः}} (\emph{prabalaḥ}) is built from \emph{pra}, \textquotedblleft{}very\textquotedblright{}, and \emph{bala}, \textquotedblleft{}strength\textquotedblright{}; MW gives \emph{prabala} as strong, mighty and violent (MW).
\end{cousinweb}

\begin{grammarlens}[title={What you've built}]
One more word for getting around.
\end{grammarlens}

\subsection*{Your turn}
\begin{itemize}
\raggedright
  \item \emph{Say it:} \emph{prabalaḥ}
  \item \emph{Say it:} \emph{prabalaḥ}, once more
  \item \emph{Say it:} \emph{samudre prabalaḥ vāyuḥ vahati.} — A violent wind blows over the sea.
  \item \emph{Recall:} say the Sanskrit for a breeze, then the Sanskrit for a tempest, then say \emph{prabalaḥ} again
\end{itemize}

\begin{tcolorbox}[breakable,colback=teal!4,colframe=teal!35!black,title={Before you move on}]
What does \sk{प्रबलः} mean? (\textquotedblleft{}Strong; violent\textquotedblright{}.) Say it once more.
\end{tcolorbox}
\section[\texorpdfstring{\sk{आन्दोलयति} (āndolayati)}{आन्दोलयति (āndolayati)}]{\sk{आन्दोलयति} (āndolayati) — sets swaying; swings}

\label{lesson:SA-C439-andolayati}



\begin{quote}
Before the new one: say the Sanskrit for a tempest, then the Sanskrit for violent.
\end{quote}

\subsection*{You'll want to know: \sk{आन्दोलयति}}
\textbf{\sk{आन्दोलयति}} — \emph{āndolayati} — \textquotedblleft{}sets swaying; swings\textquotedblright{}.

\begin{quote}
\sk{पवनः} \sk{शाखाम्} \sk{आन्दोलयति।}
\end{quote}

\begin{quote}
\emph{pavanaḥ śākhām āndolayati.}
\end{quote}

\begin{quote}
The breeze sets the branch swaying.
\end{quote}

Past (it …): \sk{पवनः} \sk{शाखाम्} \sk{आन्दोलितवान्}, \emph{pavanaḥ śākhām āndolitavān}.

\begin{cousinweb}
\textbf{\sk{आन्दोलयति}} (\emph{āndolayati}) is built on \emph{āndola}, \textquotedblleft{}swinging, a swing\textquotedblright{}, for which MW gives no derivation, so its origin is uncertain; MW gives \emph{āndolayati} as to swing or agitate, with the participle \emph{āndolita}, \textquotedblleft{}shaken, swung\textquotedblright{} (MW).
\end{cousinweb}

\begin{grammarlens}[title={What you've built}]
One more word for getting around.
\end{grammarlens}

\subsection*{Your turn}
\begin{itemize}
\raggedright
  \item \emph{Say it:} \emph{āndolayati}
  \item \emph{Say it:} \emph{āndolayati}, once more
  \item \emph{Say it:} \emph{pavanaḥ śākhām āndolayati.} — The breeze sets the branch swaying.
  \item \emph{Recall:} say the Sanskrit for a tempest, then the Sanskrit for violent, then say \emph{āndolayati} again
\end{itemize}

\begin{tcolorbox}[breakable,colback=teal!4,colframe=teal!35!black,title={Before you move on}]
What does \sk{आन्दोलयति} mean? (\textquotedblleft{}Sets swaying; swings\textquotedblright{}.) Say it once more.
\end{tcolorbox}
\section[\texorpdfstring{\sk{मन्दानिलः} (mandānilaḥ)}{मन्दानिलः (mandānilaḥ)}]{\sk{मन्दानिलः} (mandānilaḥ) — a gentle breeze}

\label{lesson:SA-C439-mandanilah}



\begin{quote}
Before the new one: say the Sanskrit for violent, then the Sanskrit for sets swaying.
\end{quote}

\subsection*{You'll want to know: \sk{मन्दानिलः}}
\textbf{\sk{मन्दानिलः}} — \emph{mandānilaḥ} — \textquotedblleft{}a gentle breeze\textquotedblright{}.

\begin{quote}
\sk{सायं} \sk{नदीतीरे} \sk{मन्दानिलः} \sk{वहति।}
\end{quote}

\begin{quote}
\emph{sāyaṃ nadītīre mandānilaḥ vahati.}
\end{quote}

\begin{quote}
In the evening a gentle breeze blows on the riverbank.
\end{quote}

A masculine noun.

\begin{cousinweb}
\textbf{\sk{मन्दानिलः}} (\emph{mandānilaḥ}) is built from \emph{manda}, \textquotedblleft{}slow, gentle\textquotedblright{}, and \emph{anila}, \textquotedblleft{}wind, air\textquotedblright{}; MW gives the compound as a gentle breeze, a zephyr (MW).
\end{cousinweb}

\begin{grammarlens}[title={What you've built}]
One more word for getting around.
\end{grammarlens}

\subsection*{Your turn}
\begin{itemize}
\raggedright
  \item \emph{Say it:} \emph{mandānilaḥ}
  \item \emph{Say it:} \emph{mandānilaḥ}, once more
  \item \emph{Say it:} \emph{sāyaṃ nadītīre mandānilaḥ vahati.} — In the evening a gentle breeze blows on the riverbank.
  \item \emph{Recall:} say the Sanskrit for violent, then the Sanskrit for sets swaying, then say \emph{mandānilaḥ} again
\end{itemize}

\begin{tcolorbox}[breakable,colback=teal!4,colframe=teal!35!black,title={Before you move on}]
What does \sk{मन्दानिलः} mean? (\textquotedblleft{}A gentle breeze\textquotedblright{}.) Say it once more.
\end{tcolorbox}
